% Generated by manual/tools/gen_reference.py from the code. Do not edit by hand.

\begin{reftabletwo}{install.sh options}{tab:ref-install-flags}{\hd{Option} & \hd{What it does}}
\texttt{-{}-dir DIR} & install directory (default \textasciitilde{}/TandemLLM) \\
\texttt{-{}-port PORT} & the API port (default 8000) \\
\texttt{-{}-host ADDR} & the address to bind (default 127.0.0.1; 0.0.0.0 serves the LAN, with no auth on the OpenAI routes) \\
\texttt{-{}-ref REF} & the engine release: tag, branch or commit (default v0.13.0) \\
\texttt{-{}-repo URL} & the git repository to clone (default the public repository) \\
\texttt{-{}-hf-\allowbreak{}token TOKEN} & a Hugging Face token (or set HF\_TOKEN; a stored 'hf auth login' works too) \\
\texttt{-{}-no-\allowbreak{}corpus} & skip the lookup corpus (the lookup drafter then copies from the conversation only) \\
\texttt{-{}-no-\allowbreak{}start} & install and configure, but do not start the engine \\
\texttt{-{}-setup} & set up in the browser instead: install the code and Wren's model only, then start setup mode (the dashboard's \#/setup page; docs/setup.md); 127.0.0.1 unless -\kern0pt-host says otherwise \\
\texttt{-{}-force-\allowbreak{}setup} & with -\kern0pt-setup, on a directory that already has a configured engine \\
\texttt{-{}-systemd} & also install a systemd user unit, tandemllm.service (off by default) \\
\texttt{-{}-yes,\allowbreak{} -\allowbreak{}y} & non-interactive: answer yes to every question \\
\texttt{-{}-dry-\allowbreak{}run} & print every step and command, change nothing \\
\texttt{-{}-uninstall} & stop the engine, remove the unit and DIR (the weights stay in the HF cache) \\
\texttt{-h,\allowbreak{} -{}-help} & this text \\
\end{reftabletwo}
\noindent Environment variables the installer reads: \texttt{TANDEM\_\allowbreak{}REF}, \texttt{TANDEM\_\allowbreak{}REPO}, \texttt{TANDEM\_\allowbreak{}DIR}, \texttt{TANDEM\_\allowbreak{}PORT}, \texttt{TANDEM\_\allowbreak{}HOST}, \texttt{HF\_\allowbreak{}TOKEN}, \texttt{HF\_\allowbreak{}HOME}.
